\scanpage{40}
\uline{Proposition} (i) Si $(LG)$ est satisfait, alors pour \deleted{tt voisinage ouvert} $\mathrm{Pic}(X)\longrightarrow\mathrm{Pic}(X_{/Y})$ est injectif, \deleted{Si} et plus précisément, on a un hom injectif
{\csname tagsleft@true\endcsname
\[
\varinjlim_{U\supset Y}\mathrm{Pic}(U)\longrightarrow\mathrm{Pic}(X_{/Y}) \tag*{(\ensuremath{\unsure{*|\alpha}})}
\]
}
avec hom de transition injectifs

(ii) Si $(LG)$ strict est satisfait, alors $(*)$ est un isom

\uline{Corollaire} Supposons que $Y$ ``est'' un diviseur de Cartier sur $X$, soit $\deleted{\unsure{\underline O|L}}\Lambda\simeq\underline L(Y)$ le faisceau inv. associé, donc $i^*(\Lambda)^{-1}\simeq J/J^2$, où $J$ est l'Idéal définissant $Y$. [$i:Y\longrightarrow X$ désigne l'injection canonique.

% slanted
\margin{Introduire aussi les $Y_n$ pour exprimer ces propriétés}

(i) Supposons $H^1(Y,\Lambda^{-n})=0$ pour $n=1,2,\ldots$, et $(LG)$. Alors $\mathrm{Pic}(X)\longrightarrow\mathrm{Pic}(Y)$ est \uline{injectif}, [plus \unsure{gén|précis} \inserted{les} \deleted{\unsure{on a un hom inj.}} \deleted{$(*)$ \unsure{est inj avec} hom} hom $\mathrm{Pic}(U)\longrightarrow\mathrm{Pic}(Y)$ sont injectifs, \deleted{\unsure{donc|d'où} bijectifs} de $=$ $\varinjlim_{U\supset Y}\mathrm{Pic}(U)\longrightarrow\mathrm{Pic}(Y)$ [et hom de transition injectifs]

(ii) Supposons $H^1(Y,\Lambda^{-n})=0$ \uuline{et} $H^2(Y,\Lambda^{-n})=0$ pour $n=1,2,\ldots$, et $(LG)$ strict. Alors $\varinjlim_{U\supset Y}\mathrm{Pic}(U)\overset{\sim}{\longrightarrow}\mathrm{Pic}(Y)$.
% archive sheet number: 39
